\documentclass[UTF8,aspectratio=169]{ctexbeamer}

\usetheme{Madrid}
\usecolortheme{seahorse}
\useinnertheme{rounded}
\setbeamertemplate{navigation symbols}{}
\setbeamertemplate{footline}[frame number]

\usepackage{booktabs}
\usepackage{array}
\usepackage{tikz}
\usetikzlibrary{arrows.meta,positioning}
\usepackage{listings}

\definecolor{okgreen}{RGB}{0,130,60}
\definecolor{warnorange}{RGB}{200,110,0}
\definecolor{mainblue}{RGB}{0,80,160}
\definecolor{codegray}{RGB}{90,90,90}

\lstset{
  basicstyle=\ttfamily\scriptsize,
  keywordstyle=\color{mainblue}\bfseries,
  commentstyle=\color{codegray}\itshape,
  stringstyle=\color{okgreen},
  frame=single, framesep=2pt, rulecolor=\color{black!30},
  backgroundcolor=\color{black!4},
  breaklines=true, columns=fullflexible, keepspaces=true,
  showstringspaces=false, aboveskip=4pt, belowskip=2pt
}

\newcommand{\okmark}{\textcolor{okgreen}{\checkmark}}
\newcommand{\warnmark}{\textcolor{warnorange}{$\triangle$}}

\title[浮点验证与规约测试进展]{浮点验证与规约测试：工作进展汇报}
\subtitle{IP 全链路验证（QCP/Rocq）与 FloatTest 位级差分测试}
\author{QCIP 验证组}
\date{2026 年 9 月}

\begin{document}

\begin{frame}
  \titlepage
\end{frame}

\begin{frame}{目录}
  \begin{columns}[T]
    \column{0.5\textwidth}
    {\scriptsize \tableofcontents[sections={1-2}]}
    \column{0.5\textwidth}
    {\scriptsize \tableofcontents[sections={3-4}]}
  \end{columns}
\end{frame}

%======================================================================
\section{背景：我们要验证什么}

\subsection{验证对象}
\begin{frame}{验证对象：星载控制软件 IP}
  \begin{itemize}
    \item \textbf{从哪里来}：卫星姿态与轨道控制系统的核心算法模块（C 语言）
      \begin{itemize}
        \item 姿态解算、陀螺数据处理、轨道计算、岁差章动、喷气推进控制\ldots
        \item 工作模式切换判断（ModeConvert 系列）
      \end{itemize}
    \item \textbf{共同特点}
      \begin{itemize}
        \item \alert{浮点密集}：大量 \texttt{float64} 四则运算、限幅、比较
        \item 频繁从\textbf{结构体/数组}中读写浮点数
        \item 调用 libm：sin/cos/sqrt，甚至 asin/atan2/exp
        \item 分支多、状态多：一个函数几十路输入输出
      \end{itemize}
    \item \textbf{形式化的必要性}：在轨软件不可维护，浮点误算代价高
  \end{itemize}
\end{frame}

\subsection{总体进展时间线}
\begin{frame}{总体进展时间线（2026-07 至 2026-08，约 7 周）}
  \centering\footnotesize
  \begin{tabular}{@{}llp{8.2cm}@{}}
    \toprule
    时间 & 路线 & 进展 \\
    \midrule
    07-02 & 验证 & 首个 IP 全链路证明完成（STSUseFlag，7 条分支全证） \\
    07-09 & 验证 & StructFilter 完成（第二代建模：float64 当 64 位整数） \\
    07-13 & 验证 & STSTransDataSave 完成（annotation 迭代 5 轮） \\
    07-14 & 验证 & \textbf{首个原生浮点验证案例}（float\_clamp / float\_store） \\
    07-15 & 验证 & WheelFriction 部分验证——浮点抽象的局限暴露 \\
    07-24 & 测试 & FloatTest 框架推广：SAM 7 题 + ModeConvert 6 题 + CS 前 2 题 \\
    08-05 & 测试 & \textbf{三角函数破题}：musl sin/cos 双侧确定性移植 \\
    08-09 & 测试 & 发现 MSVCRT sqrt 1-ulp 误舍入 $\Rightarrow$ sqrt 移植 \\
    08-10 & 测试 & 姿态 selftest 抓出 6 处共享层错误，5 题重生成 \\
    08-17 & 测试 & \textbf{iplib 24/24 全部交付}，36{,}271 条定理 \\
    08-19 & 测试 & 7 题证据链刷新 + dashboard 看板上线 \\
    \bottomrule
  \end{tabular}
  \vspace{0.8em}
  \begin{itemize}\small
    \item 另有 orbiter-new 批次 35/35 交付（35{,}035 条定理）；两条路线并行推进
    \item 07-14 原生浮点可用、08-05 三角函数解决——两类原本无法处理的程序先后纳入范围
  \end{itemize}
\end{frame}

\subsection{两条技术路线}
\begin{frame}{两条技术路线}
  \begin{columns}[T]
    \column{0.48\textwidth}
    \begin{block}{路线一：QCP 全链路形式化验证}
      \begin{itemize}\small
        \item 在 C 代码上写\textbf{规约 annotation}（前置/后置条件）
        \item 符号执行自动生成\textbf{验证条件}（VC）
        \item 在 Rocq 中人工完成\textbf{全称证明}（对所有输入成立）
        \item 结论为全称命题，但依赖工具链对浮点的支持程度
      \end{itemize}
    \end{block}
    \column{0.48\textwidth}
    \begin{block}{路线二：FloatTest 位级差分测试}
      \begin{itemize}\small
        \item 把 C 算法\textbf{逐行转写}成 Rocq 可执行规约
        \item 用 gcc 跑出真值，在采样输入上\textbf{逐比特比对}
        \item 每条测试向量 = 一条机器检查的 Rocq 定理
        \item 覆盖证明暂不可行的浮点密集程序
      \end{itemize}
    \end{block}
  \end{columns}
  \vspace{0.8em}
  \centering\small
  工具链的浮点支持程度决定程序采用哪条路线。
\end{frame}

\subsection{技术工作一览}
\begin{frame}{技术工作一览（三个方面）}
  \begin{block}{贡献 1：浮点程序的原生形式化验证机制（路线一）}
    \small 基于 Flocq 的原生 \texttt{fp32/fp64} 规约类型；用\textbf{有限性前提}在规约中显式处理 NaN/$\pm\infty$；用\textbf{类型化内存谓词} \texttt{store\_float} 表达浮点内存写入。完成 float\_clamp / float\_store 两个案例。
  \end{block}
  \begin{block}{贡献 2：确定性浮点运行时与位级规约技术（路线二）}
    \small musl sin/cos/sqrt \textbf{双侧逐行移植}，把 libm 变成确定性可复现运算；C 语言比较语义（NaN 为假）在 Rocq 侧精确复刻；混合精度、舍入模式（mode\_NE）、bit pattern 逐位对齐。
  \end{block}
  \begin{block}{贡献 3：定理化的 PBT 差分测试框架 FloatTest}
    \small ``\textbf{测试即定理}''：每条向量对应一条由 \texttt{vm\_compute} 检查的 Rocq 定理；随机 + 定向边界输入设计；\textbf{阴性自检}验证测试有效性；全链路证据固定（种子、快照、新鲜度重放）。
  \end{block}
\end{frame}

\subsection{全称证明与采样定理}
\begin{frame}{全称证明与采样定理}
  \begin{columns}[T]
    \column{0.48\textwidth}
    \begin{block}{全称证明（路线一）}
      \small
      一条定理覆盖\textbf{所有}合法输入：
      \begin{center}\footnotesize
      $\forall\, x\,\mathit{lo}\,\mathit{hi}:\; \mathit{Safe}(x,\mathit{lo},\mathit{hi}) \Rightarrow \mathit{lo} \le \texttt{clamp}(x,\mathit{lo},\mathit{hi}) \le \mathit{hi}$
      \end{center}
      证明由 Rocq 内核检查。
    \end{block}
    \column{0.48\textwidth}
    \begin{block}{测试定理（路线二）}
      \small
      每个采样点一条定理，各证一次：
      \begin{center}\footnotesize
      \texttt{t\_0861 : spec(输入$_{861}$) = 输出$_{861}$}
      \end{center}
      1000 条向量 = 1000 条定理，每条经 Rocq 内核计算验证。逐比特严格，但仅覆盖采样点。
    \end{block}
  \end{columns}
  \vspace{0.6em}
  \begin{block}{比较}
    \small 全称证明覆盖所有合法输入；采样定理仅覆盖采样点，但每条结论均由 Rocq 内核检查，严格性高于普通测试。该限制在各 case 报告中均有声明。
  \end{block}
\end{frame}

%======================================================================
\section{路线一：QCP 全链路验证}

\subsection{验证流水线}
\begin{frame}{全链路验证流水线}
  \centering
  \begin{tikzpicture}[
      stage/.style={rounded corners, draw=mainblue, thick, fill=blue!8,
                    align=center, minimum height=1.2cm, font=\small},
      arr/.style={-{Stealth[length=2.5mm]}, thick, mainblue}]
    \node[stage] (a) {C annotation\\[-1pt]\footnotesize 前置/后置条件};
    \node[stage, right=0.4cm of a] (b) {符号执行\\[-1pt]\footnotesize 自动生成 VC};
    \node[stage, right=0.4cm of b] (c) {VC 分类\\[-1pt]\footnotesize safety / return};
    \node[stage, right=0.4cm of c] (d) {Rocq 证明\\[-1pt]\footnotesize 人工完成};
    \node[stage, right=0.4cm of d] (e) {final-check\\[-1pt]\footnotesize 归档验收};
    \draw[arr] (a) -- (b);
    \draw[arr] (b) -- (c);
    \draw[arr] (c) -- (d);
    \draw[arr] (d) -- (e);
  \end{tikzpicture}
  \vspace{1em}
  \begin{itemize}\small
    \item \textbf{Annotation}：用分离逻辑描述``函数执行前内存里有什么（Require）''和``执行后保证什么（Ensure）''
    \item \textbf{Safety VC}：内存安全（不越界、类型匹配），多由自动策略完成
    \item \textbf{Return VC}：功能正确性（返回值满足后置条件），人工证明——工作量主体
    \item \textbf{final-check}：无 Admitted/Axiom 扫描、生成文件新鲜度复核、归档哈希比对
  \end{itemize}
\end{frame}

\subsection{示例 1：STSUseFlag 状态机}
\begin{frame}[fragile]{示例 1：STSUseFlag 状态机验证}
  \begin{columns}[T]
    \column{0.5\textwidth}
    \textbf{C 代码}（分支状态机，7 条路径）
    \begin{lstlisting}[language=C]
if (sSTSAtt.dataSrc == 1) {
  if (sBusSTSObj.update == 3952088175) {
    if (sSTSAtt.stsFlg == 1)
      sSTSAtt.stsFlg = 6;
    else if (sSTSAtt.stsFlg == 2)
      sSTSAtt.stsFlg = 7;
    else if (sSTSAtt.stsFlg == 4)
      sSTSAtt.stsFlg = 8;
    sSTSAtt.useOld = sSTSAtt.stsFlg;
  } else {
    if (sSTSAtt.stsFlg != 3)
      sSTSAtt.stsFlg = sSTSAtt.useOld;
  }
} else {
  sBusSTSObj.update = 0;
}
\end{lstlisting}
    \column{0.47\textwidth}
    \textbf{Annotation}（分离逻辑规约，节选）
    \begin{lstlisting}[basicstyle=\ttfamily\tiny]
/*@ Require store(&(sSTSAtt.dataSrc), data_src) *
            store(&(sBusSTSObj.update), update) *
            store(&(sSTSAtt.stsFlg), status) *
            store(&(sSTSAtt.useOld), use_old)
    Ensure exists data_src1 update1 status1 use_old1,
      STSUseFlagPost(data_src, update, status, use_old,
                     data_src1, update1, status1, use_old1)
      && store(&(sSTSAtt.stsFlg), status1) * ... */
\end{lstlisting}
    \vspace{2pt}
    {\footnotesize
    \begin{itemize}
      \item \texttt{store(p, v)}：内存地址 \texttt{p} 当前存储值 \texttt{v}
      \item \texttt{*}：分离合取——几块互不重叠的内存
      \item 后置：存在一个\textbf{新状态}，满足后状态关系 \texttt{STSUseFlagPost}
    \end{itemize}}
  \end{columns}
\end{frame}

\begin{frame}{示例 1（续）：后状态关系规约（predicate-first）}
  \begin{columns}[T]
    \column{0.55\textwidth}
    \textbf{后状态关系 \texttt{STSUseFlagPost} 的语义}（表格化）
    \begin{center}\footnotesize
    \begin{tabular}{@{}lll@{}}
      \toprule
      路径条件 & \texttt{status1} & \texttt{update1} \\
      \midrule
      数据源有效 $\wedge$ 更新标志命中 & 1$\to$6 / 2$\to$7 / 4$\to$8 & 不变 \\
      数据源有效 $\wedge$ 标志未命中 & 恢复为 \texttt{use\_old} & 不变 \\
      数据源无效 & 不变 & 清零 \\
      \bottomrule
    \end{tabular}
    \end{center}
    \vspace{4pt}
    {\small 7 条分支路径 $\Rightarrow$ 7 个 return VC，\textbf{全部证明完成}。}
    \column{0.42\textwidth}
    \begin{block}{方法论：predicate-first}
      \small
      规约描述``输入状态与输出状态之间的\textbf{数学关系}''，而非在 Rocq 中复制 C 的 if-else 控制流。
      \begin{itemize}\footnotesize
        \item 好处：规约可读、可审查、与实现解耦
        \item 反模式（归档中明确记录）：镜像 C 控制流，证明随分支爆炸
      \end{itemize}
    \end{block}
    \begin{block}{耗时构成}
      \small 证明命令本身为秒级；主要耗时在 Rocq 依赖环境编译（约 285 s/次）。
    \end{block}
  \end{columns}
\end{frame}

\subsection{浮点建模的三代演进}
\begin{frame}{浮点建模的三代演进}
  \centering
  \begin{tikzpicture}[
      box/.style={rounded corners, draw=mainblue, thick, align=center,
                  minimum height=3.1cm, text width=4.2cm, font=\small},
      arr/.style={-{Stealth[length=3mm]}, thick, mainblue}]
    \node[box] (g1) {\textbf{第一代：整数替换}\\[2pt]
      \footnotesize WheelFriction（动量轮摩擦补偿）\\
      \texttt{float64} $\to$ \texttt{unint32}\\
      公式 $\mathit{kf}\cdot h/\mathit{cut}+\mathit{vf0}\cdot\mathrm{sgn}(h)$
      被抽象为未解释函数\\[2pt]
      \textcolor{warnorange}{137 个 safety VC 无法经济证明，全部 Admitted}};
    \node[box, right=0.8cm of g1] (g2) {\textbf{第二代：不透明整数}\\[2pt]
      \footnotesize StructFilter（结构滤波 + PID）\\
      \texttt{float64} $\to$ \texttt{long long}\\
      浮点视为 64 位不透明字，仅证明整数层规格\\[2pt]
      \textcolor{okgreen}{全部 VC 证明完成，但不含浮点语义}};
    \node[box, right=0.8cm of g2] (g3) {\textbf{第三代：原生浮点}\\[2pt]
      \footnotesize float\_clamp / float\_store\\
      QCP v2.0.4 原生 \texttt{fp32/fp64}，底层 Flocq\\[2pt]
      \textcolor{okgreen}{浮点比较、有限性、内存写入全部在证明层表达}};
    \draw[arr] (g1) -- (g2);
    \draw[arr] (g2) -- (g3);
  \end{tikzpicture}
  \vspace{0.8em}
  \begin{itemize}\small
    \item 前两代为工具链不支持浮点时期的过渡方案；第三代（v2.0.4 起）开始表达浮点语义
    \item 同期对证明暂不可行的浮点密集 IP，采用 FloatTest 测试路线（后半部分）
  \end{itemize}
\end{frame}

\subsection{示例 2：float\_clamp（原生浮点）}
\begin{frame}[fragile]{示例 2：float\_clamp 验证（原生浮点）}
  \begin{columns}[T]
    \column{0.52\textwidth}
    \textbf{带 annotation 的 C 代码}（真实归档文件）
    \begin{lstlisting}[language=C,basicstyle=\ttfamily\tiny]
/*@ Extern Coq (clampFloatSafe : fp32 -> fp32 -> fp32 -> Prop) */
/*@ Extern Coq (clampFloatPost : fp32 -> fp32 -> fp32 -> fp32 -> Prop) */

float float_clamp(float x, float lo, float hi)
/*@ Require clampFloatSafe(x, lo, hi)
    Ensure clampFloatPost(x, lo, hi, __return) */
{
    if (x < lo)  return lo;
    if (x > hi)  return hi;
    return x;
}

double double_clamp(double x, double lo, double hi)
/*@ Require clampSafe(x, lo, hi)
    Ensure clampPost(x, lo, hi, __return) */
{ /* 同样三条分支 */ }
\end{lstlisting}
    \column{0.45\textwidth}
    \textbf{Rocq 侧的规约定义}（真实归档文件）
    \begin{lstlisting}[basicstyle=\ttfamily\tiny]
Definition clampFloatSafe (x lo hi : fp32) : Prop :=
  fp32_isFinite x /\ fp32_isFinite lo /\
  fp32_isFinite hi /\ fp32_le lo hi.

Definition clampFloatPost (x lo hi ret : fp32) : Prop :=
  fp32_ge ret lo /\ fp32_le ret hi.
\end{lstlisting}
    \vspace{2pt}
    {\footnotesize
    \begin{itemize}
      \item \texttt{fp32} 是 QCP v2.0.4 的\textbf{原生浮点类型}（底层 Flocq 库）
      \item 前置条件：三者\textbf{有限}且 $\mathit{lo}\le\mathit{hi}$
      \item 后置条件：$\mathit{lo}\le \mathit{ret}\le \mathit{hi}$
    \end{itemize}}
  \end{columns}
\end{frame}

\begin{frame}{示例 2（续）：NaN 与前置条件设计}
  \begin{columns}[T]
    \column{0.52\textwidth}
    \textbf{若前置条件允许 \texttt{x = NaN}}
    \begin{itemize}\small
      \item C 语言规定：NaN 参与的任何比较均为 \texttt{false}
      \item 两个分支均不成立，函数原样返回 NaN
      \item 后置条件 $\mathit{lo} \le r \le \mathit{hi}$ 对 NaN 不成立，规约不可证
    \end{itemize}
    \textbf{处理}
    \begin{itemize}\small
      \item 前置条件用 \texttt{fp32\_isFinite} 排除 NaN 与 $\pm\infty$
      \item 即显式记录假设：仅在有限输入下保证钳位性质
    \end{itemize}
    \column{0.44\textwidth}
    \begin{block}{证明要点}
      \small
      \begin{itemize}
        \item 2 个函数 $\times$ 3 条分支 = \textbf{6 个 return VC，全部证明}
        \item 规约库只有 4 个谓词 + 4 条引理
        \item 核心引理：有限浮点的比较\textbf{自反}（\texttt{x} 与 \texttt{x} 比较得 Eq）、\textbf{可交换反转}（交换两边 Lt 变 Gt）
        \item 无 Admitted，final-check 通过
      \end{itemize}
    \end{block}
    \begin{block}{说明}
      \small 本仓库首个使用原生浮点类型完成证明的案例；规约直接表达浮点比较与有限性，不再将浮点建模为整数。
    \end{block}
  \end{columns}
\end{frame}

\subsection{示例 3：float\_store}
\begin{frame}{示例 3：float\_store 与类型化内存谓词}
  \begin{columns}[T]
    \column{0.5\textwidth}
    \textbf{验证目标}
    \begin{itemize}\small
      \item \texttt{*p = x;}（\texttt{x} 为 float）在分离逻辑中对应的谓词
      \item 证明符号执行产生的写法与原生谓词等价：
      $$ p\,\#\,\text{Float} \mapsto x \;\;|-\!\!-\;\; \texttt{store\_float}\;p\;x $$
      \item 两个案例：直接写变量、写结构体字段，各 1 个 return VC，均证明完成
    \end{itemize}
    \textbf{遗留}：证明完成，但最终 canonical coqc 验收环节因环境缺失待补跑。
    \column{0.46\textwidth}
    \begin{block}{失败反例（归档于 \texttt{failing\_examples/}）}
      \small 若 annotation 中使用通用 \texttt{store(p, x)} 写浮点，符号执行报错：
      \begin{center}\footnotesize\ttfamily
        Cannot unify types fp32 and Z
      \end{center}
      \textbf{原因}：通用 \texttt{store} 仅接受 \texttt{Z} 类型值，而 \texttt{x} 为 \texttt{fp32}。浮点内存须使用类型化谓词 \texttt{store\_float}。
    \end{block}
  \end{columns}
\end{frame}

\subsection{路线一总览}
\begin{frame}{路线一总览：6 个 case 的交付状态}
  \centering\scriptsize
  \begin{tabular}{@{}llccp{4.1cm}@{}}
    \toprule
    Case & 功能 & Manual VC & 状态 & 备注 \\
    \midrule
    STSUseFlag & 星敏感器使用标志状态机 & 7 return & \okmark 完成 & 7 条分支路径全证 \\
    StructFilter & 结构滤波 + PID 控制器 & 1 return + 30 safety & \okmark 完成 & float64 建模为 64 位整数 \\
    STSTransDataSave & SRAM 248 字节分段拷贝 & 2 return & \okmark 完成 & annotation 迭代 5 轮 \\
    WheelFriction & 动量轮摩擦力矩补偿 & 4 return + 137 safety & \warnmark 部分 & 浮点抽象，safety 全 Admitted \\
    float\_clamp & fp32/fp64 钳位函数 & 2$\times$3 return & \okmark 完成 & 首个原生浮点验证 \\
    float\_store & float 内存写入谓词 & 2 return & \okmark 证明完成 & 待补最终 coqc 验收 \\
    \bottomrule
  \end{tabular}
  \vspace{1em}
  \begin{columns}[T]
    \column{0.55\textwidth}
    \begin{itemize}\small
      \item \textbf{4 个完整证明}通过 final-check、1 个部分验证、1 个证明完成待验收
      \item Annotation 是主要迭代点：STSTransDataSave 迭代 5 轮得到正确的分段规格
    \end{itemize}
    \column{0.42\textwidth}
    \begin{itemize}\small
      \item 遗留：WheelFriction 的 137 个 Admitted 等原生浮点能力重做；float\_store 待补验收
    \end{itemize}
  \end{columns}
\end{frame}

%======================================================================
\section{路线二：FloatTest 位级差分测试}

\subsection{路线二的动机}
\begin{frame}[fragile]{路线二的动机：QCP 浮点支持的限制}
  \begin{columns}[T]
    \column{0.55\textwidth}
    \textbf{目标程序的典型形态}
    \begin{lstlisting}[language=C]
// 从结构体读浮点数组 -> 浮点运算 -> 写回
tmpCp = Angle2C321(0, 0, sAtt.Psi_DA);
Cro   = MatrixMulti333(Cbiasp, tmpCp);
qri   = C2Q(Cro);          // 含三角函数与 sqrt
\end{lstlisting}
    \vspace{2pt}
    {\small\textbf{当时的工具链限制}}
    \begin{itemize}\small
      \item QCP 无法把内存中读出的 64 位数据\textbf{还原为浮点}（只当整数处理）
      \item 这批程序没有 annotation，且大量依赖 libm 三角函数
      \item 符号执行全链路无法完成
    \end{itemize}
    \column{0.41\textwidth}
    \begin{block}{FloatTest 的思路}
      \small 全称证明暂不可行时，检验一个较弱但可机器判定的性质：
      \medskip
      \begin{center}
        Rocq 规约与 C 实现在采样输入上是否\textbf{逐比特一致}
      \end{center}
      \medskip
      方法：规约逐行人工转写，与 gcc 真值逐比特比对，每个采样点形成一条 Rocq 定理。
    \end{block}
  \end{columns}
\end{frame}

\subsection{FloatTest 流程与 PBT 设计}
\begin{frame}{FloatTest 流程：测试即定理}
  \centering
  \begin{tikzpicture}[
      stage/.style={rounded corners, draw=mainblue, thick, fill=blue!8,
                    align=center, minimum height=1.5cm, text width=2.55cm, font=\small},
      arr/.style={-{Stealth[length=2.5mm]}, thick, mainblue}]
    \node[stage] (c) {原始 C 程序\\[-1pt]\footnotesize gcc -O0 真值\\[-1pt]\footnotesize 固定种子随机+定向};
    \node[stage, right=0.35cm of c] (v) {测试向量\\[-1pt]\footnotesize IEEE 754 bit pattern\\[-1pt]\footnotesize vectors.txt};
    \node[stage, right=0.35cm of v] (t) {tests.v\\[-1pt]\footnotesize 每向量一条定理\\[-1pt]\footnotesize 自动生成};
    \node[stage, right=0.35cm of t] (q) {coqc 编译\\[-1pt]\footnotesize 通过 = 测试全过\\[-1pt]\footnotesize 固定证据入口};
    \draw[arr] (c) -- (v);
    \draw[arr] (v) -- (t);
    \draw[arr] (t) -- (q);
    \node[stage, below=0.7cm of v, fill=green!8, draw=okgreen] (s) {Rocq 位级规约 spec.v\\[-1pt]\footnotesize 人工逐行转写，Flocq binary32/64};
    \draw[arr, okgreen] (s) -- (t);
  \end{tikzpicture}
  \vspace{0.6em}
  \begin{itemize}\small
    \item \textbf{真值}：原始 C 用 gcc 编译（关闭 FMA 收缩），固定种子生成随机 + 定向边界输入（阈值恰等、NaN、$\pm$Inf、$\pm 0$、饱和、计数器回绕等 7--13 类）
    \item \textbf{比对粒度}：输入输出一律按 64 位 \textbf{bit pattern} 比较，不作近似，1 比特差异即判失败
    \item \textbf{阴性自检}：每个 case 附 1 条\textbf{故意改错期望值}的定理，预期编译报错，用于验证工具链有效性
  \end{itemize}
\end{frame}

\begin{frame}{PBT 设计：六个关键决策}
  \centering\scriptsize
  \begin{tabular}{@{}lp{4.4cm}p{5.0cm}@{}}
    \toprule
    PBT 要素 & FloatTest 的设计 & 设计理由 \\
    \midrule
    被测性质 & 采样输入上 C 输出 $\equiv$ 规约输出（\textbf{逐比特}，约定 NaN==NaN） & 浮点不作近似比较，1 ulp 差异即判失败 \\
    Oracle（真值） & 原始 C + \texttt{gcc -O0}（关 FMA 收缩）+ musl 确定性移植 & 宿主 libm 不可信（sqrt 误舍入实证） \\
    输入生成 & 固定种子随机 + \textbf{7--13 类定向边界}（阈值恰等、NaN、$\pm$Inf、$\pm 0$、饱和、回绕） & 随机覆盖一般行为，定向命中浮点奇异点 \\
    检查方式 & \texttt{vm\_compute; reflexivity}——内核计算的定理 & 检查由 Rocq 内核执行，不依赖外部脚本的判断 \\
    有效性验证 & 每题 1 条\textbf{阴性自检}：故意改错，预期报错 & 变异测试思想：确认测试可检出错误 \\
    可复现性 & 固定种子 + SHA-256 快照 + 新鲜度重放 + 固定 coqc argv & 任何人任何时间重跑结论一致 \\
    \bottomrule
  \end{tabular}
  \vspace{0.5em}
  \begin{block}{与普通 PBT（如 QuickCheck）的差异}
    \small 普通 PBT 的通过记录是测试日志；FloatTest 的每条通过向量都对应一条 Rocq 定理，且规约文件可在后续证明工作中直接复用。
  \end{block}
\end{frame}

\subsection{示例 4：ModeConvert\_SBM}
\begin{frame}[fragile]{示例 4：ModeConvert\_SBM 测试流程}
  \begin{columns}[T]
    \column{0.44\textwidth}
    \textbf{第 1 步：被测 C 代码}（全部逻辑）
    \begin{lstlisting}[language=C]
// 星时超过注入时刻 + 时延 -> 转入 EIM 模式
if ((m_starTime - t0) >= dt_OrbitInject)
    m_WorkMode = WKMD_EIM;  // 0x11
// 否则工作模式不变
\end{lstlisting}
    \textbf{第 2 步：参考驱动生成向量}（\texttt{vectors.txt} 真实两行）
    \begin{lstlisting}[basicstyle=\ttfamily\tiny]
4694428428020873060 13909875606412569204 4661949024290739460 20 17
4683088663683732336 4683088663683732336 4607182418800017408 3 3
\end{lstlisting}
    {\footnotesize 列含义：starTime、t0、dt（都是 64 位 bit pattern）、旧模式字 $\to$ 新模式字}
    \column{0.53\textwidth}
    \textbf{第 3 步：人工转写的 Rocq 规约}（\texttt{spec.v} 真实全文核心）
    \begin{lstlisting}[basicstyle=\ttfamily\tiny]
Definition modeConvert_SBM_fun
    (starTime t0 dt : fp64) (workMode : Z) : Z :=
  if c_ge64 (fp64_sub starTime t0) dt
  then WKMD_EIM_Z      (* 17 = 0x11 *)
  else workMode.
\end{lstlisting}
    \textbf{第 4 步：每条向量变成一条定理}（\texttt{tests.v} 真实片段）
    \begin{lstlisting}[basicstyle=\ttfamily\tiny]
Example t_0002 :
  modeConvert_SBM_fun (f64 4683088663683732336)
    (f64 4683088663683732336) (f64 4607182418800017408) 3 = 3.
Proof. vm_compute. reflexivity. Qed.
\end{lstlisting}
    {\footnotesize \texttt{vm\_compute} 由 Rocq 内核计算左端，与右端相等则定理成立。本例 starTime = t0，差值 0 < dt，条件为假，模式保持 3。}
  \end{columns}
\end{frame}

\begin{frame}{位级测试的四个关键概念}
  \begin{columns}[T]
    \column{0.48\textwidth}
    \begin{block}{bit pattern（位模式）}
      \small 浮点数在内存中为 64 个比特；向量中的大整数即这些比特的整数值。按位比较严格于任何误差界比较。
    \end{block}
    \begin{block}{\texttt{c\_ge64}：复刻 C 的 NaN 语义}
      \small C 规定 NaN 的比较一律为假。Flocq 的原始比较遇 NaN 返回``无法比较''，因此封装 \texttt{c\_ge64}：无法比较 $\Rightarrow$ false，与 C 完全一致。
    \end{block}
    \column{0.48\textwidth}
    \begin{block}{\texttt{vm\_compute}：内核级计算}
      \small Rocq 内核的虚拟机实际执行规约函数。``通过''的含义是内核确认等式两端相等，而非外部脚本的判断。
    \end{block}
    \begin{block}{不做 extraction 的原因}
      \small 规约为人工逐行转写，转写本身可能出错，差分测试用于检出此类错误；自动提取则失去对照意义。
    \end{block}
  \end{columns}
\end{frame}

\subsection{三个技术难点}
\begin{frame}{难点 1：三角函数（15 个 case 曾无法处理）}
  \begin{columns}[T]
    \column{0.55\textwidth}
    \textbf{问题}
    \begin{itemize}\small
      \item Flocq 仅定义 IEEE 四则运算与 sqrt 的舍入语义，不含 sin/cos
      \item 各 libm 实现的三角函数结果可能不同，缺少确定的真值口径
      \item 姿态计算类 case 普遍调用三角函数，15 题受阻
    \end{itemize}
    \textbf{解决（2026-08-05）：libm 运算确定性化}
    \begin{itemize}\small
      \item 将 \textbf{musl libc 的 sin/cos}（含参数约减 \texttt{\_\_rem\_pio2}）\textbf{逐行移植}两份：
        C 侧（编译时替换系统 libm）与 Rocq 侧（\texttt{FloatTrig.v}）
      \item 两侧执行同一算法，三角函数成为\textbf{确定性、可复现}的运算
    \end{itemize}
    \column{0.41\textwidth}
    \begin{block}{移植质量自测}
      \small
      \begin{itemize}
        \item 独立 selftest：\textbf{3176 条向量逐比特一致}
        \item 与 msvcrt libm 对比：常规值 0 ulp 差，大参数约 21 ulp（musl 更准）
        \item sqrt 同样移植（musl 纯整数 Goldschmidt 迭代），自测 2073 条向量
      \end{itemize}
    \end{block}
    \begin{block}{残余边界}
      \small asin/atan2/exp 未移植：依赖它们的子函数用``\textbf{打桩 + 输入注入}''隔离——桩的输出直接作为输入注入，下游真实计算仍逐比特比对，并在每题 README 置顶声明。
    \end{block}
  \end{columns}
\end{frame}

\begin{frame}{难点 2：宿主 libm 的 sqrt 误舍入}
  \begin{columns}[T]
    \column{0.56\textwidth}
    \textbf{事件经过}（CS\_TrgtAtt\_NWM\_USU，2026-08-09）
    \begin{enumerate}\small
      \item 1025 条向量中 \textbf{t\_0861 一条}失败：输出四元数 4 个分量各差 $\pm$1 ulp
      \item 逐层排查，定位到 \texttt{sqrt}：本机 MinGW gcc 链接的老 MSVCRT 库存在罕见误舍入（x87 FSQRT 双舍入残留）
      \item 实例：$x = \texttt{0x1.de63fa81fe9bcp+0}$
        \begin{itemize}\footnotesize
          \item MSVCRT 给出 \ldots\textbf{034}（错）
          \item 正确舍入 \ldots\textbf{033}（Rocq 侧 Flocq \texttt{Bsqrt} 的结果）
        \end{itemize}
      \item 此前三个 case 约 3000 条向量未覆盖该输入，属偶然
    \end{enumerate}
    \column{0.4\textwidth}
    \begin{block}{处理与结论}
      \small 宿主 libm 不能作为真值来源，真值必须由确定性实现产生。据此移植 musl sqrt，受影响 4 题重跑全部通过。
    \end{block}
    \begin{block}{ulp 的含义}
      \small ulp（unit in the last place）是浮点末位一个比特代表的值，即浮点表示的最小差异；逐比特比对可检出该量级的错误。
    \end{block}
  \end{columns}
\end{frame}

\begin{frame}{难点 3：共享规约层的独立自测}
  \begin{columns}[T]
    \column{0.55\textwidth}
    \textbf{背景}
    \begin{itemize}\small
      \item 姿态类 case 共用一批基础算法：六种转序的 Angle2C（姿态角 $\to$ 旋转矩阵）、矩阵乘、四元数运算（QMulti / Q2C / C2Q）
      \item 提取为 C/Rocq 两侧共享库，避免每个 case 复制展开
      \item \textbf{风险}：共享层单点错误传播到所有 case
    \end{itemize}
    \textbf{独立自测（attitude selftest）}
    \begin{itemize}\small
      \item 4 个精确锚点：单位阵、绕 X/Y/Z 轴转 $\pi$
      \item 6 种转序 $\times$ 10000 条随机，从基础矩阵相乘独立重算对照
    \end{itemize}
    \column{0.41\textwidth}
    \begin{block}{发现的问题}
      \small
      \begin{itemize}
        \item \textbf{4 个 Angle2C 第一轴符号错误}
        \item \textbf{2 个 C2Q 对称项下标错误}
        \item EIM、AHM\_USU、P2P、OCM 等 5 个受影响 case 重新生成并通过
      \end{itemize}
    \end{block}
    \begin{block}{经验}
      \small 修改共享库后必须重跑所有受影响 case，避免留下``指向旧版本摘要的通过记录''。
    \end{block}
  \end{columns}
\end{frame}

\subsection{浮点处理技术汇总}
\begin{frame}{浮点处理技术汇总（问题 $\to$ 处理方式）}
  \centering\footnotesize
  \begin{tabular}{@{}p{4.6cm}p{7.6cm}@{}}
    \toprule
    浮点问题 & 处理方式 \\
    \midrule
    IEEE 舍入语义 & Flocq \texttt{Bplus/Bminus/Bmult/Bdiv}（mode\_NE 最近偶数舍入）与 \texttt{gcc -O0} 逐位对齐 \\
    内存中的浮点 & 64 位 bit pattern 注入/读出（\texttt{b64\_of\_bits}）；证明侧用类型化谓词 \texttt{store\_float} \\
    NaN 的比较语义 & 封装 \texttt{c\_ge64} 等：不可比较 $\Rightarrow$ false，精确复刻 C \\
    证明层的 NaN/$\infty$ & 前置条件 \texttt{isFinite} 显式排除，把假设写进规约 \\
    三角函数（Flocq 不支持） & musl sin/cos 双侧逐行移植（含参数约减），自测 3176 向量逐比特一致 \\
    sqrt 真值分歧 & musl sqrt（纯整数 Goldschmidt）移植 + Flocq \texttt{Bsqrt} 正确舍入 \\
    asin/atan2/exp 未移植 & 打桩 + 输入注入，边界在 README 置顶声明 \\
    混合精度 & 按 C 转换规则显式建模 binary64$\leftrightarrow$binary32 转换；\texttt{f} 后缀字面量逐行核对 \\
    \bottomrule
  \end{tabular}
  \vspace{0.6em}
  \begin{itemize}\small
    \item 上述对策基于同一套 Flocq 语义，证明与测试两条路线共用，后续做全称证明时可直接复用
  \end{itemize}
\end{frame}

\subsection{覆盖规模与证据}
\begin{frame}{测试覆盖规模与结果}
  \centering\footnotesize
  \begin{tabular}{@{}lcccl@{}}
    \toprule
    批次 & Case 数 & 向量/定理数 & 通过率 & 内容 \\
    \midrule
    \textbf{iplib} & 24/24 & 36{,}271 & 100\% & 6 个模式切换 + 18 个姿态/轨道计算 \\
    \textbf{orbiter-new} & 35/35 & 35{,}035 & 100\% & 轨道器批次，各 1000 条定向向量 \\
    \textbf{SAMCodeSynthesis} & 8/9 & 7{,}000 & 100\% & 速率/姿态控制等，最重 1 题未做 \\
    cfg\_target（早期试点） & 12 & 431 正例 & 100\% & 框架首批验证对象 \\
    \midrule
    \textbf{合计} & \textbf{79} & \textbf{$\approx$78{,}700 条定理} & \textbf{100\%} & 全部无 Admitted/Axiom \\
    \bottomrule
  \end{tabular}
  \vspace{1em}
  \begin{columns}[T]
    \column{0.55\textwidth}
    \begin{itemize}\small
      \item 典型每 case 1000 条随机 + 定向边界向量；最重 case 单题 5013 条
      \item 千条向量全链路约 10 秒；固定种子，完全可复现
    \end{itemize}
    \column{0.42\textwidth}
    \begin{itemize}\small
      \item 每题 1 条阴性自检（故意改错）全部正确报错
      \item 全链路 SHA-256 快照绑定源码/规约/向量/工具
    \end{itemize}
  \end{columns}
\end{frame}

\begin{frame}{iplib 批次：24 个 case 与证据看板}
  \begin{columns}[T]
    \column{0.55\textwidth}
    \textbf{两个程序家族}
    \begin{itemize}\small
      \item \textbf{ModeConvert}（6 题）：工作模式切换。例：EIM 题以模式字为下标读 fp64[14] 定时器数组，超时转 NWM
      \item \textbf{CS 系列}（18 题）：陀螺最小二乘/直接求逆、姿态控制喷气调度、轨道计算、IAU1976 岁差 + IAU1980 章动、点对点机动规划\ldots
      \item 最重一题 CS\_TrgtAtt\_AHM\_USU：\textbf{104 个输入 $\to$ 44 个输出}，9 段混合轨迹角速度规划
    \end{itemize}
    \textbf{共享姿态规约层}
    \begin{itemize}\small
      \item 六种 Angle2C、矩阵乘、四元数运算，C/Rocq 双侧共享
    \end{itemize}
    \column{0.41\textwidth}
    \begin{block}{证据看板（dashboard）}
      \small
      \begin{itemize}
        \item 静态网页索引：C 源码、规约声明、C$\leftrightarrow$Rocq 列映射、示例向量、复现命令
        \item 7 个 case 已刷新完整证据链（固定 coqc 检查 + 向量新鲜度逐字节比对 + 阴性控制）
        \item 其余 17 个为\textbf{历史通过、待补证据刷新}（非测试失败）
        \item 24/24 规约与测试无 \texttt{Admitted}/\texttt{Axiom}
      \end{itemize}
    \end{block}
  \end{columns}
\end{frame}

\subsection{测试发现的问题与经验}
\begin{frame}{测试发现的问题}
  \begin{enumerate}\small
    \item \textbf{宿主 libm 误舍入}：MSVCRT sqrt 存在 1-ulp 误舍入 $\Rightarrow$ 促成 musl sqrt 移植（详见难点 2）
    \item \textbf{姿态公式转写错误}：4 个 Angle2C 符号错误 + 2 个 C2Q 下标错误，被独立 selftest 抓出（详见难点 3）
    \item \textbf{驱动程序 UB}：GyroPick 的参考驱动 \texttt{printf} 格式串多写 2 个 \texttt{\%u}，打印出垃圾列，被向量生成器的\textbf{列数断言}拦截
    \item \textbf{源码笔误}：CS\_IRES\_Modify 循环误用未初始化索引 \texttt{i}（gcc 警告 + 测试边界双重记录）
    \item \textbf{混合精度陷阱}：\texttt{binary64} 实参传给 \texttt{float} 形参，要显式建模 64$\to$32$\to$64 两次转换；\texttt{0.9231f} 与 \texttt{0.9231} 精度不同，此类转写偏差可由差分测试即时检出
  \end{enumerate}
  \vspace{0.5em}
  \begin{block}{小结}
    \small 这些问题分属三类：\textbf{环境}（libm）、\textbf{规约转写}（Angle2C/C2Q）、\textbf{被测代码及其驱动}（UB、笔误）——差分测试对三类都有效。
  \end{block}
\end{frame}

\begin{frame}{测试路线：经验沉淀}
  \begin{columns}[T]
    \column{0.52\textwidth}
    \textbf{方法论}
    \begin{enumerate}\small
      \item 决定性难度因素\textbf{不是代码行数}，而是是否依赖 libm 三角/反三角/exp
      \item 真值口径固定为``原始 IP + musl 确定性移植''，结论可复现
      \item 确定性基础运算提成两侧共享层，避免复制展开导致符号/下标漂移
      \item 修改共享库后必须重跑所有受影响 case
    \end{enumerate}
    \column{0.44\textwidth}
    \textbf{工程纪律}
    \begin{enumerate}\small
      \setcounter{enumi}{4}
      \item 正式证据只接受固定 \texttt{coq\_tooling.py check} 在隔离 workspace 的固定 argv 结果
      \item 每题 README 置顶声明重建约定、打桩边界与未覆盖范围
      \item 全链路 SHA-256 快照绑定：C 源码、规约、驱动、向量、发射器
    \end{enumerate}
  \end{columns}
  \vspace{0.8em}
  \begin{block}{定位}
    \small FloatTest 是 QCP 全链路的\textbf{补充}而非替代：在证明暂时不可达的浮点程序上，用采样定理给出机器检查的一致性结论；规约库可在将来做全称证明时复用。
  \end{block}
\end{frame}

%======================================================================
\section{总结与展望}

\subsection{工作量量化}
\begin{frame}{工作量量化}
  \begin{columns}[T]
    \column{0.48\textwidth}
    \begin{block}{路线一：验证（6 case，约 2 周）}
      \small
      \begin{itemize}
        \item 人工 Rocq 证明 \textbf{$\approx$1{,}330 行}（proof\_manual 合计）
        \item 人工数学规约库 \textbf{$\approx$630 行}
        \item 人工证明 return VC \textbf{22 个} + safety VC 30 个（含拆分约 90 条引理）
        \item 最大单 case：StructFilter 生成 goal 9{,}749 行
        \item annotation 最多迭代 \textbf{5 轮}（STSTransDataSave）
      \end{itemize}
    \end{block}
    \column{0.48\textwidth}
    \begin{block}{路线二：测试（79 case，约 5 周）}
      \small
      \begin{itemize}
        \item 人工逐行转写规约 \textbf{$\approx$4{,}300 行}（三批次 spec.v + 共享库）
        \item 生成并编译通过定理 \textbf{$\approx$78{,}700 条}（tests.v 合计 $\approx$47 万行）
        \item 测试向量 \textbf{78{,}271 条}，每题 7--13 类定向边界
        \item musl sin/cos/sqrt \textbf{双侧移植 3 套}；独立自测向量 6 万+ 条
        \item 中文 README 边界声明 $\approx$2{,}600 行；静态 dashboard 1 套
      \end{itemize}
    \end{block}
  \end{columns}
  \vspace{0.4em}
  \centering\small
  合计：\textbf{85 个 case} 交付，证明与测试两条线共覆盖 IP 库的主要浮点程序。
\end{frame}

\subsection{总结}
\begin{frame}{总结：三方面技术工作，两条路线}
  \begin{columns}[T]
    \column{0.55\textwidth}
    \textbf{技术贡献}
    \begin{enumerate}\small
      \item \textbf{浮点原生验证机制}：fp32/fp64 规约 + 有限性前提 + 类型化内存谓词，规约直接表达浮点语义
      \item \textbf{确定性浮点运行时与位级规约}：musl 移植 + C 语义复刻，解决``浮点没有可信真值''的根本难题
      \item \textbf{定理化 PBT 框架 FloatTest}：测试即定理 + 阴性自检 + 证据固定，采样测试具有机器检查的结论
    \end{enumerate}
    \column{0.41\textwidth}
    \textbf{完成情况}
    \begin{itemize}\small
      \item 路线一：6 case，4 个完整证明，三代浮点建模演进
      \item 路线二：79 case，约 \textbf{78{,}700 条定理 100\% 通过}，零 Admitted/Axiom
      \item 实证发现的问题：宿主 libm 误舍入、规约转写错误、驱动 UB
    \end{itemize}
  \end{columns}
  \vspace{0.6em}
  \centering\small
  证明可达的程序做全称证明；证明暂不可达的浮点程序，用逐比特的定理化测试覆盖。
\end{frame}

\subsection{未来工作}
\begin{frame}{未来工作}
  \begin{columns}[T]
    \column{0.48\textwidth}
    \textbf{验证路线}
    \begin{itemize}\small
      \item 用原生浮点能力重做 WheelFriction 等早期抽象 case，补齐 137 个 Admitted
      \item float\_store 补跑 canonical coqc 验收
      \item 沉淀外部函数规格库；循环不变式自动生成
    \end{itemize}
    \column{0.48\textwidth}
    \textbf{测试路线}
    \begin{itemize}\small
      \item 补齐 17 个 archived case 的最新证据链
      \item 完成遗留的最重 case（SAM GyroAttiDetermine：矩阵求逆 + fp64 floor）
      \item asin/atan2/exp 的确定性移植，消除打桩边界
      \item \textbf{测试 $\to$ 证明衔接}：已验证的 Flocq 规约库可直接复用于全称性定理（如限幅输出界）
    \end{itemize}
  \end{columns}
  \vspace{1em}
  \centering\small
  两条路线的衔接点：FloatTest 规约库 + QCP 原生浮点支持，为浮点 IP 的全称证明提供基础。
\end{frame}

\begin{frame}
  \centering
  {\Huge 谢谢！}\\[1.5em]
  {\large 欢迎讨论与指正}
\end{frame}



\end{document}
